\paragraph{Decoupled effect and permission readouts.}
For every workspace file $f$ represented at a recovered write site, Tideline evaluates two lightweight neural readouts over the frozen backbone features (each a 128-unit GELU MLP): an \emph{effect readout} predicting whether the candidate performs no modeled write, an \texttt{append}, or a \texttt{replace}, and a separate \emph{permission readout} outputting independent binary grants for \texttt{may\_append} and \texttt{may\_replace}. Explicit task clauses and statically resolved write sites refine these readouts into per-file sets $E_f$ and $P_f$, and the gate denies execution whenever $\bigcup_f (E_f \setminus P_f) \neq \varnothing$.

Factoring authorization across both files and write modes prevents two common failure modes. First, a file-name match alone cannot distinguish an instruction to append an audit entry from permission to overwrite the file in place. Second, verifying that a primary report write is authorized cannot discharge an unauthorized write to a secondary log or ledger. Because the two readouts are small and decoupled from the frozen 8B backbone, external policy engines can also override $P_f$ directly while feedback updates adapt only the readouts (Section~\ref{sec:continual-design}). Routing frequencies across the contextual verification checks are detailed in Appendix~\ref{app:construction}.
